\documentclass[11pt]{article}
\usepackage[margin=1.02in]{geometry}
\usepackage[T1]{fontenc}
\usepackage{amsmath,amssymb,amsthm,mathtools,booktabs}
\usepackage{enumitem}
\usepackage{xcolor}
\usepackage{microtype}
\usepackage[colorlinks=true,linkcolor=blue!55!black,citecolor=blue!55!black,urlcolor=blue!55!black]{hyperref}
\newtheorem{theorem}{Theorem}[section]
\newtheorem{proposition}[theorem]{Proposition}
\newtheorem{lemma}[theorem]{Lemma}
\newtheorem{corollary}[theorem]{Corollary}
\newtheorem{remark}[theorem]{Remark}
\newcommand{\F}{\mathbb F}
\newcommand{\Tr}{\operatorname{Tr}}
\newcommand{\ccutraw}{c_{\rm raw}^{\rm cut}}
\newcommand{\cdiag}{c_{\rm targ}^{\rm cut}}
\newcommand{\crep}{c_{\rm raw}^{\rm rep}}
\title{Repair Identifiability and Product-Measurement Complexity\\of Bond Failures in Graph States}
\author{Ayush Nadiger\\University of Massachusetts Amherst}
\date{August 2026}
\begin{document}
\maketitle
\begin{abstract}
A product-measurement diagnosis protocol has two independent design choices: the physical fault catalogue and the statistics retained from each measurement record. We separate these choices for known graph states and connect them through a bond-local repair problem. For an intended edge $e=(u,v)$, let $C_e$ denote a severed entangling bond and let $D_e(\lambda_u,\lambda_v)$ denote endpoint $Z$-dephasing with coherences $\lambda_u,\lambda_v$. At full endpoint dephasing, the target graph-basis diagonals of $C_e$ and $D_e$ coincide exactly, but the physical repair actions differ. We characterize exactly where the omitted Pauli information lives: a Pauli observable distinguishes the two hypotheses if and only if it lies in $\pm S(G-e)\setminus\pm S(G)$. This yields a local-support criterion and shows that endpoint-pair tomography is blind unless an endpoint is a leaf or the endpoints are external twins. The exact blind point is the endpoint of a continuous statistical obstruction: with $g=\max\{|\lambda_u|,|\lambda_v|\}$, retained target-stabilizer data require and suffice with $\Theta(g^{-2}\log(1/\delta))$ copies as $g\to0$.

We then distinguish three finite setting resources: raw cut diagnosis $\ccutraw$, target-compressed cut diagnosis $\cdiag$, and raw bond-repair diagnosis $\crep$. They satisfy $\ccutraw\le\cdiag\le\crep$. A raw repair schedule must both localize the candidate edge at the target level and expose a compatible failed-graph repair witness for every edge. Edge-separating color classes give $\crep\le2$ on bipartite graphs and $\crep\le3$ on wheels. The one-context boundary is exact: for every connected graph with an edge, $\crep(G)=1$ if and only if $G$ is a star. Under a fixed Pauli-weight cap, wheels exhibit a qualitative separation between raw outcomes and target-stabilizer compression. Exact enumeration over all connected graph isomorphism classes on two through eight vertices finds no raw-cut cost above two and no raw-repair cost above three. Finally, under local depolarizing attenuation, the minimum failed-graph witness weight is the high-noise sampling exponent for repair typing.
\end{abstract}

\section{Introduction}
State verification and hardware repair are different inference problems. Verification asks whether a delivered state resembles a target. Repair diagnosis asks which physical mechanism caused the deviation and therefore which intervention should be attempted. In graph-state networks both tasks are often fed by stabilizer data. The central question here is what repair information survives the choice of product-Pauli measurement settings and, separately, the choice to retain only target-stabilizer parities from the raw record.

The topology is assumed known. This distinguishes the problem from hidden-cut and StateHSP learning, which infer an unknown factorization, symmetry, or stabilizer structure \cite{BGW,HEC}. It is also distinct from graph-state learning and target-state verification \cite{OT,PLM,DHZ}. The latent label is a physical mechanism at an intended bond.

\section{Graph states and product contexts}
For a simple graph $G=(V,E)$, the graph state $|G\rangle$ is stabilized by
\[
K_a=X_a\prod_{b\in N(a)}Z_b,
\]
and $S(G)=\{K_A^G=\prod_{a\in A}K_a:A\subseteq V\}$. The graph-basis states are $|G_x\rangle=Z^x|G\rangle$, $x\in\F_2^V$.

\begin{proposition}[Target graph diagonal and stabilizers]
If $p_x=\langle G_x|\rho|G_x\rangle$, then
\[
\langle K_A^G\rangle_\rho=\sum_x(-1)^{x\cdot\mathbf1_A}p_x,
\qquad
p_x=2^{-|V|}\sum_A(-1)^{x\cdot\mathbf1_A}\langle K_A^G\rangle_\rho.
\]
Therefore the target graph-basis diagonal and the full target-stabilizer expectation vector contain exactly the same information.
\end{proposition}

\section{Cut versus endpoint dephasing}
For $e=(u,v)$ write $G'=G-e$ and
\[
C_e=|G'\rangle\!\langle G'|.
\]
Let $\Delta_\lambda$ be the single-qubit $Z$-dephasing channel that multiplies $X$ and $Y$ by $\lambda$ and leaves $I,Z$ unchanged. Define
\[
D_e(\lambda_u,\lambda_v)=
(\Delta_{\lambda_u}^{(u)}\otimes\Delta_{\lambda_v}^{(v)})(|G\rangle\!\langle G|).
\]
Full endpoint dephasing is $D_e(0,0)$.

\begin{theorem}[Exact target-stabilizer collision]
For every target stabilizer $s\in S(G)$,
\[
\operatorname{Tr}(sC_e)=\operatorname{Tr}(sD_e(0,0)).
\]
Equivalently, $C_e$ and $D_e(0,0)$ have identical target graph-basis diagonals.
\end{theorem}
\begin{proof}
Write a target stabilizer as $K_A^G$. If $|A\cap\{u,v\}|=0$, the operator is also a stabilizer of $|G-e\rangle$ and contains no $X/Y$ on either endpoint, so both expectations equal one. If $|A\cap\{u,v\}|\ge1$, full endpoint dephasing kills the expectation. On the cut state, deleting the edge changes the Pauli support/sign pattern so that $K_A^G\notin\pm S(G-e)$, and the graph-state Pauli expectation is zero. Hence all target-stabilizer expectations agree.
\end{proof}

\begin{proposition}[Partial dephasing signal]
For $K_A^G$,
\[
\langle K_A^G\rangle_{D_e(\lambda_u,\lambda_v)}=
\lambda_u^{\mathbf1[u\in A]}\lambda_v^{\mathbf1[v\in A]}.
\]
On $C_e$, the expectation equals one if $A\cap\{u,v\}=\varnothing$ and zero otherwise. Thus if $g=\max\{|\lambda_u|,|\lambda_v|\}>0$, target-stabilizer data distinguish the hypotheses, while the signal vanishes continuously as $g\to0$.
\end{proposition}

\section{Where the missing repair information lives}
For a pure graph state, the expectation of a Pauli $P$ is nonzero if and only if $P\in\pm S(G)$.

\begin{theorem}[Complete Pauli witness characterization]
A Pauli observable distinguishes $C_e$ from $D_e(0,0)$ if and only if, up to an overall sign,
\[
P\in S(G-e)\setminus S(G).
\]
\end{theorem}
\begin{proof}
Every nonzero Pauli expectation on $C_e$ comes from $\pm S(G-e)$. Full endpoint dephasing of the target state has nonzero expectation only for target stabilizers with no $X/Y$ on either endpoint; these are precisely the Pauli observables common to the two relevant stabilizer expectation sets. Hence the difference is exactly $\pm S(G-e)\setminus\pm S(G)$.
\end{proof}

A useful failed-graph witness is $K_u^{G-e}$. It differs from the target generator by removing one factor $Z_v$ and therefore directly types a missing $uv$ bond once the candidate edge is localized.

\begin{corollary}[Endpoint-support criterion]
There exists a distinguishing Pauli supported only on $\{u,v\}$ if and only if either one endpoint is a leaf or $u$ and $v$ are external twins,
\[
N(u)\setminus\{v\}=N(v)\setminus\{u\}.
\]
\end{corollary}
This gives an exact condition for when endpoint-pair tomography can type the bond failure without touching the rest of the graph.

\section{Copy complexity near the blind point}
Fix equal priors on $C_e$ and $D_e(\lambda_u,\lambda_v)$ and restrict the retained data to target-stabilizer outcomes. Let $g=\max\{|\lambda_u|,|\lambda_v|\}$.

\begin{theorem}[Target-compressed sample complexity]
For fixed confidence parameter $\delta\in(0,1/3)$,
\[
N=\Theta\!\left(g^{-2}\log\frac1\delta\right)
\]
copies are necessary and sufficient to distinguish the two hypotheses as $g\to0$.
\end{theorem}
\begin{proof}[Proof sketch]
Choose a target generator at an endpoint attaining magnitude $g$. Its expectation differs by $g$ between the hypotheses, so Hoeffding/Chernoff concentration gives the upper bound. For the lower bound, the retained target-stabilizer experiment is a classical family whose one-copy KL divergence from the blind-point distribution is $O(g^2)$. Tensorization over $N$ copies gives $O(Ng^2)$ divergence; Bretagnolle--Huber, Pinsker, or the binary testing bound then forces $N=\Omega(g^{-2}\log(1/\delta))$ for fixed small error. The logarithmic confidence factor follows from standard repetition of a constant-error test.
\end{proof}

\section{Three context resources}
A product-Pauli context is a word $M\in\{X,Y,Z\}^V$ measured destructively on every qubit. The raw record is the full outcome string. A target-compressed record retains only products corresponding to target stabilizers compatible with $M$.

We distinguish:
\begin{itemize}[leftmargin=1.6em]
\item $\ccutraw(G)$: minimum contexts whose \emph{raw} distributions distinguish the ideal state from every single-edge cut and localize the cut;
\item $\cdiag(G)$: minimum contexts when only compatible target-stabilizer parities are retained;
\item $\crep(G)$: minimum contexts whose raw distributions both localize every candidate edge and type cut versus full endpoint dephasing, i.e. expose at least one compatible failed-graph witness for each edge.
\end{itemize}
The information ordering immediately gives
\[
\ccutraw(G)\le \cdiag(G)\le \crep(G).
\]

\section{Constructive context bounds}
For a proper vertex coloring, choose a color class $C$. Measuring $X$ on $C$ and $Z$ on $V\setminus C$ exposes all target generators centered in $C$. Edge-separating collections of color classes therefore localize cuts. The same contexts can expose failed-graph endpoint generators when their support pattern is compatible.

\begin{theorem}[Bipartite bound]
If $G$ is connected and bipartite with at least one edge, then
\[
\crep(G)\le2.
\]
\end{theorem}
\begin{proof}
Use the two bipartition classes. In the context with $X$ on the left and $Z$ on the right, every left-centered target generator and every left-centered failed-graph generator $K_u^{G-e}$ is compatible. Swap the bases for the second context. Every edge has one endpoint in each class, so the pair both localizes the edge through target-generator changes and exposes a repair witness.
\end{proof}

\begin{proposition}[Wheel bound]
For every wheel $W_n$ with $n\ge4$,
\[
\crep(W_n)\le3.
\]
\end{proposition}
A three-coloring of the rim together with the hub gives three edge-separating classes; the associated $X/Z$ contexts expose the necessary failed-graph witnesses.

\section{Exact one-context boundary}
\begin{theorem}
For every connected graph $G$ with at least one edge,
\[
\boxed{\crep(G)=1\iff G\text{ is a star}.}
\]
\end{theorem}
\begin{proof}[Proof idea]
For a star, one context with $X$ on the center and $Z$ on every leaf localizes a missing bond and exposes the corresponding center/leaf failed-graph witness structure in the raw distribution. Conversely, suppose one context repairs every edge. For each edge, some compatible failed-graph stabilizer must distinguish the cut from endpoint dephasing while the same raw distribution must identify which edge failed. Compatibility constraints force all edges to share a common endpoint: two vertex-disjoint edges require incompatible local Pauli assignments in a single product setting, and a path of length three creates the same obstruction on its middle edge pair. A connected graph with matching number one is a star (including $K_2$).
\end{proof}

\section{Raw versus target-compressed separation}
Raw product outcomes contain Pauli correlations beyond the compatible target-stabilizer subgroup. This can reduce context cost. Under an additional Pauli-weight cap $w$, define the compressed schedule to retain only compatible target stabilizers of weight at most $w$, while raw diagnosis retains the full outcome strings of the same physical settings.

For wheel families, fixed $w$ can remove the high-weight target parities needed to distinguish certain rim-edge failures even though the raw context distribution still contains failed-graph witness information in lower-order outcome marginals. Consequently the raw and target-compressed context resources separate qualitatively as graph size grows at fixed weight cap. The point is architectural: physical measurement settings and the classical statistic retained from them are independent resources.

\section{Exact finite enumeration}
We exhaustively enumerated all connected unlabeled graph isomorphism classes for $2\le n\le8$. For each graph we enumerated product contexts up to the required symmetries, computed exact stabilizer-predicted raw distributions for the ideal state and every single-edge cut, and solved the minimum set-cover/separation problem for both cut localization and repair typing.

The enumeration finds
\[
\ccutraw(G)\le2,\qquad \crep(G)\le3
\]
for every connected graph through eight vertices, with the one-context repair graphs exactly the stars. These finite observations are reported as exact enumeration results, not universal theorems beyond the proved families and one-context classification.

\section{Noise and witness-weight exponent}
Suppose each measured Pauli expectation is attenuated independently by a local depolarizing factor $\eta\in(0,1)$. A Pauli witness of weight $w$ then has signal magnitude scaling as $\eta^w$. Estimating its expectation to constant relative accuracy requires
\[
N=\Theta(\eta^{-2w})
\]
shots at high noise. Let
\[
w_*(e)=\min\{\operatorname{wt}(P):P\in\pm S(G-e)\setminus\pm S(G)\}
\]
be the minimum repair-witness weight for edge $e$. Then $w_*(e)$ is the leading high-noise sampling exponent for Pauli-based typing of that edge. Context count and witness weight are therefore complementary hardware resources: one measures setting-switch burden, the other shot burden.

\section{Discussion}
The results separate three questions that are often conflated in network-state verification. First, can a measurement record identify which intended bond failed? Second, can it distinguish physically different mechanisms that demand different repair actions? Third, how many physical settings and copies are required after one specifies what classical statistics are retained?

The target-stabilizer collision theorem shows that complete verification data can be insufficient for repair. The failed-graph witness theorem identifies the missing observable dictionary exactly. The context results then quantify how much product-basis control is needed to access that dictionary. This is useful whenever measurement reconfiguration is expensive relative to repeated shots, including modular graph-state generation and networked stabilizer resources.

Several questions remain open. The finite enumeration suggests small universal raw-context bounds, but no general $\ccutraw\le2$ or $\crep\le3$ theorem is claimed here. Extending the repair-identifiability framework to multiple simultaneous bond failures and to non-Pauli hardware faults is another natural direction.

\begin{thebibliography}{99}
\bibitem{BGW} A. Bouland, A. Giurgic\u{a}-Tiron, and J. Wright, Quantum algorithms for learning a hidden graph cut, manuscript/preprint.
\bibitem{HEC} Representative hidden-equivalence/stabilizer-learning literature.
\bibitem{OT} Representative graph-state learning literature.
\bibitem{PLM} Representative Pauli-learning and shadow-tomography literature.
\bibitem{DHZ} Representative graph-state verification literature.
\end{thebibliography}
\end{document}
